\section{Conclusion}
\label{sec:conclusion}
Score-debiased kernel density estimation (SD-KDE) offers improved asymptotic accuracy over classical KDE, but its reliance on an empirical score has historically made it substantially slower in practice. 
This paper shows that the dominant cost of empirical SD-KDE is not inherently prohibitive: by re-ordering the computation to expose matrix-multiplication structure, both KDE evaluation and the score numerator can be expressed as Tensor Core--accelerable GEMMs plus lightweight elementwise operations, while avoiding materializing full pairwise interaction matrices via streaming accumulation. 
In the 16-D setting, this hardware-aligned formulation yields large end-to-end gains---up to $47\times$ faster than a strong Torch SD-KDE baseline and $3{,}300\times$ faster than scikit-learn KDE---and scales to previously infeasible regimes, completing SD-KDE on $\sim$1M training points and $\sim$131k queries in 2.3\,s on a single GPU. 

 Overall, our results suggest that practical nonparametric estimators can often be unlocked by hardware-aware reformulations that maximize GPU utilization. 
 Future directions include extending the approach beyond $d$ multiples of 16, supporting richer bandwidth parameterizations and kernels, further reducing non-Tensor-Core overheads (e.g., exponentials and atomics), and developing nonnegativity-preserving approximations and multi-GPU variants.
